\documentclass[11pt,a4paper]{article}

\usepackage[margin=2.5cm]{geometry}
\usepackage{amsmath,amssymb}
\usepackage{enumitem}
\usepackage[most]{tcolorbox}

% ---------- Notation ----------
\newcommand{\vb}[1]{\mathbf{#1}}
\newcommand{\vect}[1]{
    \begin{pmatrix}#1
\end{pmatrix}}
\newcommand{\PP}{\mathbb{P}}
\newcommand{\RR}{\mathbb{R}}

% ---------- Consistent layout ----------
% Every exercise:      \exercise{n}
% Every sub-question:  \task{label}{short description}
% Every final result:  \begin{answer} ... \end{answer}
\newcommand{\exercise}[1]{
\section*{Exercise #1}}
\newcommand{\task}[2]{
\subsection*{#1\quad\normalfont\textit{#2}}}
\newtcolorbox{answer}{colback=gray!8,colframe=gray!60,boxrule=0.5pt,
    arc=2pt,left=6pt,right=6pt,top=4pt,bottom=4pt,
fontupper=\normalsize,before upper={\textbf{Answer.}\ }}

\setlength{\parindent}{0pt}
\setlength{\parskip}{6pt}

\title{Assignment 2 -- Theoretical Exercises}
\author{Computer Vision, 2026 \\ Arvid Kolstad}
\date{}

\begin{document}
\maketitle

% =====================================================================
\exercise{1}

\task{}{Projective ambiguity for uncalibrated cameras}

Start from the camera equation and insert the identity $T^{-1}T = I$,
where $T$ is any
projective transformation of 3D space (an invertible $4\times 4$ matrix):
\[
    \lambda\vb{x} = P\vb{X} = P\left(T^{-1}T\right)\vb{X}.
\]
If we set
\[
    \tilde{P} = PT^{-1} \qquad\text{and}\qquad \tilde{\vb{X}} = T\vb{X},
\]
we have
\[
    \lambda\vb{x} = P\vb{X} = \tilde{P}\tilde{\vb{X}}.
\]
Since $T$ is invertible, $\tilde{P}$ is still a $3\times 4$ matrix of
rank 3, i.e.\ a valid camera.
With unknown calibration, any such camera is allowed:
$\tilde{P} = PT^{-1} = K[R\;\vb{t}]\,T^{-1}$. To recover $K$ we would
use an RQ decomposition
of $P = [KR\;\;K\vb{t}]$, decomposing $A = KR$. But since $P$ itself
is unknown, this cannot be done.
\begin{answer}
    $(\tilde{P}, \tilde{\vb{X}}) = (PT^{-1},\, T\vb{X})$ gives
    exactly the same image projections as
    $(P, \vb{X})$ for every projective transformation $T$. Hence the
    reconstruction can only be
    recovered up to an unknown projective transformation of 3D space.
\end{answer}

% =====================================================================
\exercise{2}

\task{}{Calibrated cameras and the remaining ambiguity}

For a calibrated camera, with a normalized coordinates, we still
have a rotation matrix and a translation vector that is unknown.
We can also scale it how we want without the equation breaking. Therefore
we reduce the amounts of T:s that allowed to rotations,
translation and scaling. So there is still ambiguity, but it's less
than before. To be explicit T is a 4x4 matrix that for the
constraints of excercise 1 has $4 \dot 4 -1 =
15$ dimensions of freedom (because of scaling of the matrix makes no
difference). With these constraints the matrix has 3 from rotations,
3 from translations and 1 from scaling which equals 7 DOFs.

\begin{answer}
    The ambiguity goes down from 15 DOFs to 7 DOFs.
\end{answer}

% =====================================================================
\exercise{3}

\[
    K = \vect{f&0&x_0\\0&f&y_0\\0&0&1}.
\]

\task{(a)}{Verify $K^{-1}$ and the factorization $K^{-1} = AB$}

Use $K^{-1} = \frac{1}{\det K}\operatorname{adj}(K)$, with $\det K =
f^2$. The cofactor matrix is
\[
    \operatorname{adj}(K)^T = \vect{f&0&0\\0&f&0\\-fx_0&-fy_0&f^2}
    \;\Longrightarrow\;
    \operatorname{adj}(K) = \vect{f&0&-fx_0\\0&f&-fy_0\\0&0&f^2},
\]
so
\[
    K^{-1} =
    \vect{\tfrac{1}{f}&0&-\tfrac{x_0}{f}\\[2pt]0&\tfrac{1}{f}&-\tfrac{y_0}{f}\\[2pt]0&0&1}.
\]
Multiplying the factors:
\[
    AB =
    \vect{\tfrac{1}{f}&0&0\\[2pt]0&\tfrac{1}{f}&0\\[2pt]0&0&1}\vect{1&0&-x_0\\0&1&-y_0\\0&0&1}
    =
    \vect{\tfrac{1}{f}&0&-\tfrac{x_0}{f}\\[2pt]0&\tfrac{1}{f}&-\tfrac{y_0}{f}\\[2pt]0&0&1}.
\]
\begin{answer}
    Both expressions equal $K^{-1}$, so the inverse and the
    factorization check out.
\end{answer}

\task{(b)}{Geometric interpretation of $A$ and $B$}

$A$ only scales the lengths of the current coordinates, by $1/f$ in
$x$, $1/f$ in $y$ and $1$ in $z$.
$B$ is a shearing in the $x$- and $y$-direction: the larger the
$z$-coordinate, the larger the shift in
the $x$- and $y$-direction.
\begin{answer}
    $A$ is a scaling by $1/f$. $B$ shears the $x$- and
    $y$-coordinates in proportion to $z$.
\end{answer}

\task{(c)}{Interpretation of $K^{-1}$}

By multiplying with $K^{-1}$ we remove the inner parameters, so that
the remaining camera only
depends on rotation and translation, i.e.\ the camera's position and
orientation.
Because of this, the principal point ends up in the origin:
\[
    K^{-1}\vect{x_0\\y_0\\1} = \vect{x_0/f - x_0/f\\ y_0/f -
    y_0/f\\ 1} = \vect{0\\0\\1}.
\]
A point at distance $f$ from the principal point ends up at distance
$1$ after normalization:
\[
    K^{-1}\vect{x_0+f\\y_0\\1} = \vect{(x_0+f)/f - x_0/f\\ y_0/f -
    y_0/f\\ 1} = \vect{1\\0\\1}.
\]
\begin{answer}
    $K^{-1}$ removes the inner parameters. The principal point is
    mapped to the origin $(0,0)$, and
    a point at distance $f$ from the principal point is mapped to a
    point at distance $1$ from the origin.
\end{answer}

\task{(d)}{Normalize $(0,300)$ and $(800,300)$}

With $f = 400$, $x_0 = 400$ and $y_0 = 300$, the formula from (a) gives
\[
    K^{-1} =
    \vect{\tfrac{1}{400}&0&-1\\[2pt]0&\tfrac{1}{400}&-\tfrac{3}{4}\\[2pt]0&0&1}.
\]
\[
    K^{-1}\vect{0\\300\\1} = \vect{-1\\0\\1},\qquad
    K^{-1}\vect{800\\300\\1} = \vect{1\\0\\1}.
\]
\begin{answer}
    $(0,300) \mapsto (-1,\,0)$, \quad $(800,300) \mapsto (1,\,0)$.
\end{answer}

\task{(e)}{Angle between the viewing rays}

The viewing rays have the directions of the normalized points,
$\vb{v}_1 = (-1,0,1)^T$ and $\vb{v}_2 = (1,0,1)^T$. From linear algebra,
\[
    \vb{v}_1^T\vb{v}_2 = \|\vb{v}_1\|\,\|\vb{v}_2\|\cos\alpha
    \quad\text{and}\quad
    \vb{v}_1^T\vb{v}_2 = -1 + 0 + 1 = 0
    \;\Longrightarrow\; \cos\alpha = 0.
\]
\begin{answer}
    $\alpha = 90^\circ$.
\end{answer}

\task{(f)}{Same camera center and principal axis}

\textbf{Camera center.} Let $\vb{X}'$ be in the null space of
$[R\;\vb{t}]$, i.e.\ $[R\;\vb{t}]\vb{X}' = \vb{0}$. Then
\[
    K[R\;\vb{t}]\vb{X}' = K\vb{0} = \vb{0},
\]
so $\vb{X}'$ also solves $K[R\;\vb{t}]\vb{X} = \vb{0}$. Conversely,
since $K$ is invertible,
$K[R\;\vb{t}]\vb{X}' = \vb{0}$ implies $[R\;\vb{t}]\vb{X}' =
K^{-1}\vb{0} = \vb{0}$.
The two cameras therefore have the same null space, i.e.\ the same
camera center.

\textbf{Principal axis.} The last row of $K$ is $(0,\,0,\,1)$, so the
last row of
$K[R\;\vb{t}]$ is the last row of $[R\;\vb{t}]$. With $\vb{t} =
(t_1,t_2,t_3)^T$, the camera equation
$\lambda\vb{x} = K[R\;\vb{t}]\vb{X}$ gives
\[
    \lambda = R_{31}X + R_{32}Y + R_{33}Z + t_3 .
\]
$\lambda$ is the depth in these coordinates, so this equation
describes planes orthogonal to the
principal axis, and $(R_{31}, R_{32}, R_{33})$ is the direction of
the camera. This does not depend
on $K$ in any way.
\begin{answer}
    $K[R\;\vb{t}]$ and $[R\;\vb{t}]$ have the same null space (camera
    center) and the same last row
    $(R_{31}, R_{32}, R_{33})$ (principal axis).
\end{answer}

% =====================================================================
\exercise{4}

\task{(a)}{Verify the product $KR$}

\[
    KR = \vect{a&b&c\\0&d&e\\0&0&f}\vect{R_1^T\\R_2^T\\R_3^T}
    = \vect{aR_1^T + bR_2^T + cR_3^T\\ dR_2^T + eR_3^T\\ fR_3^T}.
\]

\task{(b)}{$R_3$ and $f$}

\[
    P =
    \vect{800\sqrt2&0&2400\sqrt2&8000\\-700\sqrt2&2800&700\sqrt2&9800\\-\sqrt2&0&\sqrt2&6}
\]
\[
    \Longrightarrow\quad
    KR =
    \vect{800\sqrt2&0&2400\sqrt2\\-700\sqrt2&2800&700\sqrt2\\-\sqrt2&0&\sqrt2}.
\]
The last row gives
\[
    f\,(R_{31},\,R_{32},\,R_{33}) = (-\sqrt2,\,0,\,\sqrt2).
\]
Since $\|R_3\| = 1$ and $f > 0$, we get $f = \|(-\sqrt2, 0, \sqrt2)\| = 2$.
\begin{answer}
    $f = 2$, \quad $R_3^T =
    \left(-\tfrac{1}{\sqrt2},\,0,\,\tfrac{1}{\sqrt2}\right)$.
\end{answer}

\task{(c)}{$R_2$, $d$ and $e$}

The second row gives $dR_2^T + eR_3^T = (-700\sqrt2,\,2800,\,700\sqrt2)$, with
$eR_3^T = e\left(-\tfrac{1}{\sqrt2},0,\tfrac{1}{\sqrt2}\right)$:
\[
    \begin{cases}
        -700\sqrt2 = -\dfrac{e}{\sqrt2} + dR_{21}\\[6pt]
        2800 = dR_{22}\\[2pt]
        700\sqrt2 = \dfrac{e}{\sqrt2} + dR_{23}
    \end{cases}
\]
Since $R_3 \perp R_2$, we have $R_3^TR_2 = \tfrac{1}{\sqrt2}(-R_{21}
+ R_{23}) = 0$, so $R_{21} = R_{23}$.
Adding the first and third equations gives $0 = 2dR_{21}$, so $R_{21}
= R_{23} = 0$ (since $d > 0$).
Then $\|R_2\| = 1$ gives $R_{22} = 1$, and therefore $d = 2800$ and $e = 1400$.
\begin{answer}
    $R_2^T = (0,\,1,\,0)$, \quad $d = 2800$, \quad $e = 1400$.
\end{answer}

\task{(d)}{$R_1$, $a$, $b$ and $c$}

Because $R_1 \perp R_2$ and $R_1 \perp R_3$ (with $\det R = 1$),
$R_1^T = \left(\tfrac{1}{\sqrt2},\,0,\,\tfrac{1}{\sqrt2}\right)$. The
first row then gives
\[
    \begin{cases}
        800\sqrt2 = \dfrac{a}{\sqrt2} - \dfrac{c}{\sqrt2}\\[6pt]
        0 = b\\[2pt]
        2400\sqrt2 = \dfrac{a}{\sqrt2} + \dfrac{c}{\sqrt2}
    \end{cases}
    \;\Longrightarrow\;
    \begin{cases}
        a - c = 1600\\
        b = 0\\
        a + c = 4800
    \end{cases}
    \;\Longrightarrow\;
    a = 3200,\; b = 0,\; c = 1600.
\]
\begin{answer}
    $R_1^T = \left(\tfrac{1}{\sqrt2},\,0,\,\tfrac{1}{\sqrt2}\right)$,
    \quad $a = 3200$, \quad $b = 0$, \quad $c = 1600$.
\end{answer}

\task{(e)}{Focal length, skew, aspect ratio and principal point}

\[
    R =
    \vect{\tfrac{1}{\sqrt2}&0&\tfrac{1}{\sqrt2}\\[2pt]0&1&0\\[2pt]-\tfrac{1}{\sqrt2}&0&\tfrac{1}{\sqrt2}},
    \qquad
    K = \vect{3200&0&1600\\0&2800&1400\\0&0&2}
    \sim \vect{1600&0&800\\0&1400&700\\0&0&1}.
\]
$R$ is a rotation by $45^\circ$ around the $y$-axis. Comparing with
$K = \vect{\gamma f & s & x_0\\ 0 & f & y_0\\ 0&0&1}$:
\begin{answer}
    Focal length $f = 1400$, \quad skew $s = 0$, \quad aspect ratio
    $\gamma = \tfrac{1600}{1400} = \tfrac{8}{7}$,
    \quad principal point $(x_0, y_0) = (800,\,700)$.
\end{answer}

% =====================================================================
\exercise{5}

\task{(a)}{$\min_{\vb{v}}\|M\vb{v}\|^2$ always has the minimum value $0$}

Since $\|M\vb{v}\|^2 \geq 0$, the smallest possible value is $0$.
With no constraints on $\vb{v}$
we can set $\vb{v} = \vb{0}$ and always get $\|M\vb{0}\|^2 = 0$.
\begin{answer}
    The minimum is always $0$, attained at $\vb{v} = \vb{0}$.
\end{answer}

\task{(b)}{$\|M\vb{v}\|^2 = \|\Sigma V^T\vb{v}\|^2$ and $\|V^T\vb{v}\| = 1$}

With $M = U\Sigma V^T$ and $U^TU = I$ (since $U$ is orthogonal):
\begin{align*}
    \|M\vb{v}\|^2 &= (M\vb{v})^T(M\vb{v}) = (U\Sigma
    V^T\vb{v})^T(U\Sigma V^T\vb{v})
    = \vb{v}^TV\Sigma^T\underbrace{U^TU}_{I}\Sigma V^T\vb{v}\\
    &= (\Sigma V^T\vb{v})^T(\Sigma V^T\vb{v}) = \|\Sigma V^T\vb{v}\|^2.
\end{align*}
If $\|\vb{v}\|^2 = 1$, the same trick with $VV^T = I$ (since $V$ is
orthogonal) gives
\[
    \|V^T\vb{v}\| = \sqrt{\vb{v}^TVV^T\vb{v}} = \sqrt{\vb{v}^T\vb{v}}
    = \|\vb{v}\| = 1,
\]
where we take the positive root since a length cannot be negative.
\begin{answer}
    $\|M\vb{v}\|^2 = \|\Sigma V^T\vb{v}\|^2$, and $\|V^T\vb{v}\| =
    \|\vb{v}\| = 1$.
\end{answer}

\task{(c)}{The equivalent problem
$\min_{\|\tilde{\vb{v}}\|^2=1}\|\Sigma\tilde{\vb{v}}\|^2$}

$\Sigma$ has the singular values $\sigma_1 \geq \sigma_2 \geq \dots
\geq 0$ on its diagonal.
Since $V^T$ is orthogonal, we can see $\tilde{\vb{v}} = V^T\vb{v}$ as
a change of coordinates which
preserves length, so $\|\tilde{\vb{v}}\| = 1$ exactly when
$\|\vb{v}\| = 1$. Together with
$\|M\vb{v}\|^2 = \|\Sigma\tilde{\vb{v}}\|^2$, the two problems take
the same values, so they have the
same minimal value. To get the actual vector back, we transform
$\tilde{\vb{v}}^* \to \vb{v}^*$ with
$V$, i.e.\ $\vb{v}^* = V\tilde{\vb{v}}^*$.

We can see that there must also be at least two different solutions
to this problem when considering that the negative version of the
same vector must have the same answer because of the fact that the
norm of a vector scaled by -1 is the same as.
\begin{answer}
    (12) and (9) have the same minimal value, and $\vb{v}^* =
    V\tilde{\vb{v}}^*$. Concider the scaling -1 of the vector to
    always get a second solution.
\end{answer}

\task{(d)}{The last column of $V$ solves (9)}

Because the singular values in $\Sigma$ are in decreasing order,
$\tilde{\vb{v}} = (0,\dots,0,1)^T$
is a solution to (12). This means that
\[
    \vb{v}^* = V\vect{0\\ \vdots\\ 0\\ 1} = \text{the last column of } V,
\]
which is therefore our explicit solution to (9).
we use the projection

\begin{answer}
    The last column of $V$ is a solution to (9).
\end{answer}

% =====================================================================
\exercise{6}

\task{}{Camera for the original problem}

\[
    \tilde{\vb{x}} \sim N\vb{x}
    \;\Longrightarrow\;
    \vb{x} \sim N^{-1}\tilde{\vb{x}} \sim N^{-1}\tilde{P}\vb{X},
\]
and comparing with $\vb{x} \sim P\vb{X}$:
\begin{answer}
    $P = N^{-1}\tilde{P}$.
\end{answer}

\end{document}
